\bagiandepan{PETUNJUK PENGISIAN}

\begingroup\itshape
\noindent Halaman ini dan halaman Peta Kewajiban dihapus dari naskah final. Isi setiap bab dan subbab dalam bentuk paragraf. Teks miring berawalan ``Petunjuk:'' menjelaskan isi yang diharapkan dan harus dihapus setelah bagian itu ditulis; teks abu-abu dalam kurung siku diganti dengan isi Anda.\par
\vspace{4pt}
\noindent Tabel hanya dipakai untuk menyajikan data yang memang berbentuk tabel, diberi nomor dan judul (Tabel 2.1 di atas tabel), dan dirujuk dalam narasi. Jangan menulis isi bab atau subbab di dalam tabel atau kotak.\par
\endgroup
\vspace{6pt}

\begin{itemize}[leftmargin=0.53cm,itemsep=3pt,topsep=0pt,parsep=0pt]
\item Judul penelitian maksimal 16 kata dan ditulis sama persis di sampul, lembar pengesahan, pernyataan, dan presentasi.
\item Pengerjaan bersifat individu. Topik boleh beririsan dengan payung riset dosen, tetapi kontribusi orisinal Anda harus dapat dibedakan (Pedoman TA 3.1.2).
\item Naskah seminar proposal berisi Bab I sampai III; laporan sidang akhir berisi Bab I sampai V lengkap. Lihat Peta Kewajiban.
\item Ukuran naskah UNESCO (lebar 15,5 cm $\times$ tinggi 23 cm), Times New Roman 11 pt, spasi 1,5 (spasi 1 untuk abstrak, isi tabel, dan daftar pustaka), rata kiri-kanan (justify). Margin atas 2 cm, kiri 2,5 cm, bawah 2 cm, kanan 2 cm; dicetak bolak-balik dengan margin cermin (margin 2,5 cm selalu di sisi jilid). Judul tabel di atas tabel dan judul gambar di bawah gambar.
\item Sitasi gaya IEEE; minimal 15 sumber, minimal 70\% terbit dalam 3 tahun terakhir; prioritaskan jurnal dan prosiding bereputasi (Pedoman TA 8.7).
\item Kode sumber dikelola di git dengan commit berkala (minimal mingguan); laporan dikerjakan dengan riwayat versi. Akses diberikan ke pembimbing sejak awal dan ke penguji saat sidang (Pedoman TA 4.2.2).
\item Code walkthrough WAJIB pada sidang akhir bila TA memuat implementasi perangkat lunak. Siapkan Peta Kode pada Lampiran 4.
\item Similarity maksimal 25\% dan kutipan langsung maksimal 10\% (laporan Turnitin lengkap); Halaman Pernyataan AI wajib diisi (Pedoman TA 5.7, 4.2.1).
\item Bimbingan minimal 4 kali per semester, dicatat di Form Kendali Bimbingan atau Kartu Kendali TASYS.
\end{itemize}

\clearpage
\bagiandepan{PETA KEWAJIBAN PER TAHAP}

\noindent{\singlespacing\itshape Seminar proposal (TA 1): dilaksanakan sebelum penelitian; yang dinilai adalah kelayakan rencana penelitian (Bab I sampai III). Sidang akhir (TA 2): dilaksanakan setelah penelitian selesai; yang dinilai adalah hasil, pembahasan, dan pemahaman mahasiswa.\par}

\vspace{4pt}
\noindent{\textbf{\textcolor{wajibred}{WAJIB}}: harus ada. \ \textbf{\textcolor{wajibred}{WAJIB bila berlaku}}: bila kondisinya terpenuhi. \ \textbf{\textcolor{opsionalolive}{Opsional}}: boleh dikosongkan. \ \textcolor{placeholdergray}{Tidak diperlukan}: tidak diisi.\par}

\vspace{6pt}
\begingroup\singlespacing\footnotesize\noindent
\begin{tabularx}{\textwidth}{|p{5.14cm}|>{\centering\arraybackslash}X|>{\centering\arraybackslash}X|}
\hline
\rowcolor{shadebiru}\centering\textbf{Bagian} & \textbf{Seminar proposal} & \textbf{Sidang akhir} \\\hline
Halaman judul, pengesahan, orisinalitas & \cellcolor{shadewajib}\WAJIB & \cellcolor{shadewajib}\WAJIB \\\hline
Pernyataan penggunaan AI & \cellcolor{shadeheader}\TIDAKPERLU & \cellcolor{shadewajib}\WAJIB \\\hline
Kata pengantar & \cellcolor{shadeopsional}\OPSIONAL & \cellcolor{shadewajib}\WAJIB \\\hline
Abstract dan Abstrak & \cellcolor{shadeopsional}\OPSIONAL & \cellcolor{shadewajib}\WAJIB \\\hline
BAB I Pendahuluan & \cellcolor{shadewajib}\WAJIB & \cellcolor{shadewajib}\WAJIBF \\\hline
BAB II Tinjauan Pustaka & \cellcolor{shadewajib}\WAJIB & \cellcolor{shadewajib}\WAJIBF \\\hline
BAB III Metodologi & \cellcolor{shadewajib}\WAJIB & \cellcolor{shadewajib}\WAJIBF \\\hline
BAB IV Hasil dan Pembahasan & \cellcolor{shadeheader}\TIDAKPERLU & \cellcolor{shadewajib}\WAJIB \\\hline
BAB V Kesimpulan dan Saran & \cellcolor{shadeheader}\TIDAKPERLU & \cellcolor{shadewajib}\WAJIB \\\hline
Daftar pustaka & \cellcolor{shadewajib}\WAJIB & \cellcolor{shadewajib}\WAJIB \\\hline
Lampiran 1 (Form Kendali Bimbingan) & \cellcolor{shadewajib}\WAJIB & \cellcolor{shadewajib}\WAJIB \\\hline
Lampiran 2 (berita acara seminar proposal) & \cellcolor{shadeheader}\TIDAKPERLU & \cellcolor{shadewajib}\WAJIB \\\hline
Lampiran 3 sampai 5 (kode, peta kode, data) & \cellcolor{shadeopsional}\OPSIONAL & \cellcolor{shadewajib}\WAJIB \\\hline
Lampiran 6 sampai 8 (similarity, log AI, riwayat versi) & \cellcolor{shadeheader}\TIDAKPERLU & \cellcolor{shadewajib}\WAJIB \\\hline
Lampiran 9 (daftar riwayat hidup) & \cellcolor{shadeheader}\TIDAKPERLU & \cellcolor{shadeopsional}\OPSIONAL \\\hline
\end{tabularx}
\endgroup
\clearpage
